\documentclass{article}
\usepackage{amsmath,amssymb,amsthm,hyperref}
\newtheorem{proposition}{Proposition}
\begin{document}
\title{Binary hyperplanes: the joint correction equation}
\maketitle

\paragraph{Status.}
This note gives exact written identities and a scoped obstruction to a
particular shared-jet implementation.  It is not a Lean theorem or an
improvement of the original worst-case head bound.  All table columns
and all point memories below are independently arbitrary.

\section{The larger cell and its interference}

Let $G=\mathbb F_2^r$, let $\chi_\alpha(x)=(-1)^{\alpha\cdot x}$,
and include the zero direction in the following sums.  For each
direction $\alpha\in G$ choose a pole $p_\alpha$, and for each point
$g\in G$ and column $j$ choose a token $t_{gj}<\min_\alpha p_\alpha$.
The hyperplane memories are
\[
 Z_{\alpha,e}(j)=\sum_{g:\alpha\cdot g=e}
             \frac{1}{p_\alpha-t_{gj}}.
\]
For a query bit $b$, choose fixed residues $\eta_b(\alpha)$ and put
\[
 F_b(t)=\sum_\alpha\frac{\eta_b(\alpha)}{p_\alpha-t},
 \qquad
 K_b(z,t)=\sum_\alpha\frac{\chi_\alpha(z)\eta_b(\alpha)}
                                      {p_\alpha-t}.
\]
The formal main score obtained by selecting the hyperplane through $x$
is exactly
\[
 \sum_\alpha\eta_b(\alpha)Z_{\alpha,\alpha\cdot x}(j)
 =\frac12\sum_g F_b(t_{gj})
  +\frac12\sum_g K_b(x+g,t_{gj}).                 \tag{1}
\]
Indeed, the equality indicator is
$(1+\chi_\alpha(x+g))/2$; distributing the finite sums proves (1).
The diagonal term has $K_b(0,t_{xj})=F_b(t_{xj})$.  Therefore the
exact correction needed to obtain that diagonal lookup is
\[
 \frac12\sum_g F_b(t_{gj})
 +\frac12\sum_{g\ne x}K_b(x+g,t_{gj})
 -\frac12 F_b(t_{xj}).                           \tag{2}
\]
The first term is independent of the row and can be absorbed in a
column correction.  The second is not: for every independently chosen
$g$-memory it retains a whole Walsh transform of the Cauchy spectrum.

In the affine plane, a distinct pair of points determines one line
direction, so its corresponding correction is one Cauchy fraction.
Here a distinct pair lies on $2^{r-1}$ of the selected hyperplanes.
Replacing their sum by one Cauchy fraction is an additional, generally
false interpolation equation.  It does not follow from shared geometry.

For example, with $p_\alpha=D-\sum_i w_i\alpha_i$ and
$\eta_b(\alpha)=1$, a character involving every coordinate gives
\[
 K_b((1,\ldots,1),t)
 =\sum_{A\subseteq[r]}\frac{(-1)^{|A|}}
                             {D-t-\sum_{i\in A}w_i}.
\]
This is the mixed finite-difference kernel already studied in
\texttt{JointCellFilteredKernels.lean}.  It has many poles when the
subset sums are distinct.  The binary geometry has therefore exposed
the earlier tensor correction problem, not removed it.

\section{What shared polynomial repair does and does not save}

A parity on a support $S$ can be recovered from a suitable degree
$|S|$ jet of a reciprocal affine statistic, after subtracting lower
Boolean degrees.  In a joint table compiler, those subtractions have
coefficients depending on the column tokens.  If residues depend on
the queried bit, the subtraction coefficients depend on both $b$ and
$j$.  They cannot be replaced by one label-independent correction.

On the full $r$-cube, the union of all proper-degree row monomials has
dimension $2^r-1$.  Paying one column-memory correction per row
monomial and bit therefore costs $a(2^r-1)$ heads up to the constant
cost of the admissible monomial primitive.  This cancels the proposed
$a$-fold gain from $2^r$ point memories.  This is a cost statement for
that implementation, not a lower bound for all rational scores.

There is a second collapse to avoid.  If a proposed cancellation turns
all column dependence into just $M$ arbitrary real residues against
fixed row--bit features, each column lies in a fixed linear space of
dimension at most $M$.  The number of its strict sign vectors on $aM$
queries is at most $2\sum_{k=0}^{M-1}\binom{aM-1}{k}$, which is
$(O(a))^M$.  It cannot represent all $2^{aM}$ sign vectors for growing
$a$.  Keeping genuine, independently variable column poles may evade
this particular linear-space obstruction.

\section{The remaining moment-cancellation question}

Orthogonal-array identities could in principle cancel lower jets
across many hyperplanes before decoding.  Such an identity must hold
with the actual token-dependent coefficients, not merely with uniform
weights.  For instance,
\[
 \sum_\alpha\chi_S(\alpha)=0
 \quad\not\Longrightarrow\quad
 \sum_\alpha\frac{\chi_S(\alpha)}{(p_\alpha-t_{\alpha j})^k}=0
\]
for independent $t_{\alpha j}$.  Fixing every token except one and
varying that remaining token already disproves the implication.
An admissible jointly chosen token dictionary satisfying the needed
moment equations while retaining arbitrary table signs remains open.

\end{document}
